Het \texttt{ss} commando is de vervanger van \texttt{netstat}. Met dit commando kunnen we zien welke poorten er op ons systeem gebruikt worden. De meeste opties van \texttt{netstat} werken ook bij \texttt{ss}, wat handig is voor degenen die al vertrouwd zijn met \texttt{netstat}.

De belangrijkste opties voor \texttt{ss} zijn:
\begin{description}
\item[-n] vervang IP-adressen niet door namen, dit scheelt een hoop snelheid
\begin{lstlisting}[language=bash]
$ ss -n
\end{lstlisting}

\item[-t] Laat alle TCP poorten zien
\begin{lstlisting}[language=bash]
$ ss -t
\end{lstlisting}

\item[-u] Laat alle UDP poorten zien
\begin{lstlisting}[language=bash]
$ ss -u
\end{lstlisting}

\item[-l] Laat alle Listening poorten zien
\begin{lstlisting}[language=bash]
$ ss -l
\end{lstlisting}

\item[-p] Laat alle gebruikte poorten zien met het erbij behorende proces. Dit kunnen we alleen zien als de root-gebruiker
\begin{lstlisting}[language=bash]
$ sudo ss -p
\end{lstlisting}

\end{description}

Natuurlijk kunnen we opties combineren, zo kunnen we alle luisterende poorten laten zien op TCP, met hun bijbehorende processen zonder resolving door aan \texttt{ss} de volgende opties meet te geven:
\begin{lstlisting}[language=bash]
$ sudo ss -tlnp
\end{lstlisting}

